\section{Conclusion}

When state-of-the-art language models excel on one trustworthiness benchmark, they systematically excel on others---a pattern we set out to explain. Through large-scale meta-analysis of 4,561 models, we demonstrate that this correlation is not a scale confound to be dismissed but a signal revealing a latent factor: Generalized Representational Coherence.

GRC persists after controlling for model size and training recency ($\lambda_1 = 2.277$, 60\% variance, $p = 0.001$). It correlates with behavioral stability ($\rho = 0.405$), linking the factor to a plausible mechanism. Instruction-tuning increases both stability and factor scores with large effect sizes ($d \approx 2.0$).

The pervasive $\rho = 0.80$--$0.87$ correlation across trustworthiness benchmarks is not noise. It is evidence that beneath the surface-level diversity of evaluation tasks lies a shared structure---one that instruction-tuning improves and that future work can target directly.
